\chapter{Service IA et Scouting}

\section{Vue d'ensemble}

Le service de scouting IA est un microservice FastAPI dedie aux fonctionnalites d'intelligence artificielle et de scouting sportif. Il complete l'ecosysteme en fournissant:

\begin{itemize}[leftmargin=*]
  \item \textbf{ML01}: Score de potentiel sportif avec explications des facteurs
  \item \textbf{ML02}: Analyse de l'evolution sportive (progression, stabilite, regression)
  \item \textbf{ML03}: Prediction du churn (risque d'abandon, raisons probables, actions recommandees)
  \item \textbf{SC01}: Recherche avancee, comparaison des joueurs, generation de shortlists
\end{itemize}

\section{Stack technique}

\begin{table}[H]
\centering
\begin{tabular}{p{0.35\textwidth} p{0.55\textwidth}}
\toprule
\textbf{Composant} & \textbf{Technologie} \\
\midrule
Framework API & FastAPI \\
Langage & Python 3.9+ \\
Base de donnees & SQLAlchemy + SQLite/PostgreSQL \\
ML/Data Science & scikit-learn, pandas, numpy \\
API Documentation & Swagger UI (auto-generee) \\
Tests & pytest + pytest-asyncio \\
\bottomrule
\end{tabular}
\caption{Synthese stack service IA}
\end{table}

\section{Architecture du service}

\subsection*{Structure modulaire}

\begin{lstlisting}[style=codeblock]
app/
|-- main.py              # Point d'entree FastAPI
|-- api/
|   |-- routes/
|   |   |-- health.py    # Health check
|   |   |-- data.py      # Data ingestion endpoints
|   |   |-- ml.py        # ML scoring endpoints
|   |   |-- scouter.py  # Scouting features
|-- services/
|   |-- ml_service.py    # Logique ML
|   |-- feature_engineering.py
|   |-- sync_service.py # Sync avec football-academy
|-- db/
|   |-- models.py        # Modeles SQLAlchemy
|   |-- session.py      # Gestion session DB
|-- tests/
|   |-- test_scoring.py  # Tests unitaires ML
\end{lstlisting}

\section{Fonctionnalites ML}

\subsection*{ML01 - Score de potentiel sportif}

L'endpoint \texttt{GET /api/v1/ml/potential/\{player\_external\_id\}} calcule un score de potentiel base sur:

\begin{itemize}[leftmargin=*]
  \item Statistiques de performance (buts, passes decisives, notes moyennes)
  \item Assiduite aux entrainements
  \item Historique des blessures
  \item Age et position
\end{itemize}

\textbf{Reponse API:}
\begin{lstlisting}[style=codeblock]
{
  "player_external_id": 101,
  "potential_score": 78.5,
  "factors": {
    "performance": 82.0,
    "attendance": 75.0,
    "health": 80.0,
    "age_factor": 76.0
  },
  "explanation": "Fort potentiel offensif avec bonne assiduite"
}
\end{lstlisting}

\subsection*{ML02 - Analyse de l'evolution}

L'endpoint \texttt{GET /api/v1/ml/evolution/\{player\_external\_id\}} analyse la tendance sur une fenetre temporelle configurable:

\begin{itemize}[leftmargin=*]
  \item \textbf{Progression}: Amelioration significative des performances
  \item \textbf{Stabilite}: Performances constantes
  \item \textbf{Regression}: Baisse des performances
\end{itemize}

\subsection*{ML03 - Prediction du churn}

L'endpoint \texttt{GET /api/v1/ml/churn/\{player\_external\_id\}} estime le risque d'abandon:

\begin{lstlisting}[style=codeblock]
{
  "player_external_id": 101,
  "churn_probability": 0.35,
  "risk_level": "medium",
  "likely_reasons": [
    "payment_delays",
    "decreasing_attendance"
  ],
  "recommended_actions": [
    "contact_parent",
    "offer_payment_plan",
    "schedule_meeting"
  ]
}
\end{lstlisting}

\section{Fonctionnalites Scouting}

\subsection*{SC01 - Recherche avancee}

L'endpoint \texttt{GET /api/v1/scouter/players/search} permet une recherche multi-criteres:

\begin{itemize}[leftmargin=*]
  \item Filtrage par position, age, division
  \item Seuils sur potentiel minimum et churn maximum
  \item Tri par differentes metriques
\end{itemize}

\textbf{Exemple de requete:}
\begin{lstlisting}[style=codeblock]
GET /api/v1/scouter/players/search?
    position=MID&
    min_potential=65&
    max_churn=0.4&
    limit=20
\end{lstlisting}

\subsection*{Generation de shortlists}

L'endpoint \texttt{POST /api/v1/scouter/shortlists/generate} cree des listes de joueurs selon differentes strategies:

\begin{itemize}[leftmargin=*]
  \item \textbf{balanced}: Equilibre entre potentiel et faible risque
  \item \textbf{potential}: Priorite au potentiel pur
  \item \textbf{safe}: Priorite aux joueurs a faible risque
\end{itemize}

\section{Integration avec football-academy}

\subsection*{Synchronisation des donnees}

L'endpoint \texttt{POST /api/v1/data/sync/football-academy} synchronise automatiquement:

\begin{itemize}[leftmargin=*]
  \item Donnees joueurs depuis \texttt{/football-academy/api/players}
  \item Donnees paiements depuis \texttt{/football-academy/api/payments}
\end{itemize}

\subsection*{Configuration de l'integration}

\begin{lstlisting}[style=codeblock]
# .env
FOOTBALL_BACKEND_BASE_URL=http://localhost:8091
SYNC_INTERVAL_HOURS=24
\end{lstlisting}

\section{Tests et qualite}

\subsection*{Tests unitaires}

Les tests couvrent:
\begin{itemize}[leftmargin=*]
  \item Calcul des scores de potentiel
  \item Analyse d'evolution
  \item Prediction de churn
  \item Recherche et filtrage
\end{itemize}

\subsection*{Validation des modeles}

Les modeles ML sont valides via:
\begin{itemize}[leftmargin=*]
  \item Cross-validation
  \item Mesures de performance (Accuracy, F1, MAE, RMSE)
  \item Analyse des features importantes
\end{itemize}

\section{Deploiement}

\subsection*{Dockerisation}

Le service est conteneurise avec Docker:

\begin{lstlisting}[style=codeblock]
FROM python:3.9-slim
WORKDIR /app
COPY requirements.txt .
RUN pip install -r requirements.txt
COPY . .
CMD ["uvicorn", "app.main:app", "--host", "0.0.0.0", "--port", "8010"]
\end{lstlisting}

\subsection*{Integration Docker Compose}

Le service s'integre dans l'ecosysteme via docker-compose.yml:

\begin{lstlisting}[style=codeblock]
services:
  scouting-ai:
    build: ./scouting-ai-service
    ports:
      - "8010:8010"
    environment:
      - FOOTBALL_BACKEND_BASE_URL=http://football-academy:8091
    depends_on:
      - football-academy
\end{lstlisting}